%% kapa: exam=1 round=21 year=2010 period=2 form=A subject=accounting subject_ko=회계학 seq=39 number=39
%% source: Q-Net 감정평가사 기출문제 (cst00309) artlSeq=5024181 file=extracted/1008/5024181/2010년 제21회 감정평가사 1차시험 기출문제/2교시.hwp
%% conversion: kordoc 4.24.51 -> tools/gen1.py (auto); review=proofread-sonnet; flags=none
\begin{question}{39}
(주)한국은 20×0년에 선박제조업을 개시하였으며, 20×0년 생산 및 원가 자료는 아래와 같다. (주)한국은 직접작업시간을 기준으로 제조간접원가를 예정배부하고 있다. 20×0년도 발생한 실제 제조간접원가는 총 ￦400,000이고, 예정 제조간접원가는 ￦320,000이며, 예정 직접작업시간은 4,000시간이었다. LNG선은 완성된 즉시 선주에게 인도되었으며, 유람선은 20×0년 말까지 인도되지 않았다. 제조간접원가 배부차이를 총원가기준으로 배분할 때 LNG선의 매출원가는 얼마인가?

{\small\centering
\begin{tabular}{|l|r|r|r|r|}
\hline
 & \multicolumn{1}{c|}{LNG선} & \multicolumn{1}{c|}{유람선} & \multicolumn{1}{c|}{화물선} & \multicolumn{1}{c|}{합계} \\
\hline
직접재료원가(원) & 340,000 & 244,000 & 126,000 & 710,000 \\
직접노무원가(원) & 180,000 & 220,000 & 130,000 & 530,000 \\
실제 직접작업시간(시간) & 1,500 & 1,200 & 1,800 & 4,500 \\
완성도(\%) & 100 & 100 & 70 & - \\
\hline
\end{tabular}\par}

\begin{choices}
  \item ￦410,000
  \item ￦480,000
  \item ￦532,000
  \item ￦574,000
  \item ￦656,000
\end{choices}
\end{question}
\answer{5}
